\section{The commutant and structural consequences}\label{sec:structure}
\begin{proof}[Proof of Theorem~\ref{thm:direct-model} and Corollary~\ref{cor:direct-algebra}]
Proposition~\ref{prop:fibers} proves the fiber classification and norm
statement in Theorem~\ref{thm:direct-model}, and
Proposition~\ref{prop:normquasi} gives the uniform power bound and
nonzeroness. Proposition~\ref{prop:V} constructs the commuting
contraction with its dense tail images. Section~\ref{sec:subspaces}
proves transitivity of the algebra generated by $S$, $V$, and the
scalar multipliers. These establish all assertions of the theorem.

That transitive algebra is contained in $\{S\}'$, so any closed
$\{S\}'$-invariant subspace is zero or all of $H_0$. The nonzero
norm-quasinilpotent operator $S$ is nonscalar. The commutant lemma
\cite[\CommutantImportLocator]{CircleCompanion}, with this $S$ on the
complex Hilbert space $H_0$, now implies that $\{S\}'$ is a proper
strongly closed unital complex algebra. Together with the lattice
conclusion, it is transitive.

Finally, the finite residue classes form a countable generating algebra
on $X$, so $L^2(X)$ and $H_0$ are separable. The space $H_0$ is
infinite-dimensional. Every infinite-dimensional separable complex
Hilbert space is unitarily isomorphic to $H_0$; conjugation by such a
unitary transfers $S$, its norm estimates, its commutant, and the
invariant-subspace conclusions. This proves the corollary.
\end{proof}

\subsection{Ordinary invariant subspaces and compactness}
For every $h\in\Z$, the subspace $L^2(X;\Ktail{h})$ is closed,
nonzero, proper, and $S$-invariant. It is not invariant under $V$:
the image of the constant field $e_h$ has its nonzero coordinate at
$h-1$ almost everywhere. Thus the construction separates ordinary
invariance from hyperinvariance in an explicit way.

Each fiber $S_x$ is compact. Indeed $a_k^\pm\ge c(k+1)$ implies
$\beta_j(x)\to0$ as $|j|\to\infty$, and cutting off all but finitely
many shift weights approximates $S_x$ in operator norm. The operator
$S$ on $H_0$ is not compact. There is an infinite orthonormal sequence
$(\varphi_n)$ in $L^2(X)$, since the $N_m$ disjoint residue classes give
$N_m$ orthogonal indicators for every $m$. The orthonormal vectors
$\varphi_n\otimes e_0$ have images
\[
 S(\varphi_n\otimes e_0)=e^{-c}\varphi_n\otimes e_1.
\]
These images are orthogonal with a common positive norm and have no
norm-convergent subsequence.

\subsection{The multiplier algebra has infinite multiplicity}
Let $Q_0$ be the projection onto $L^2(X)\otimes\C e_0$. This gives a
concrete witness to the properness of the commutant: on $1\otimes e_0$,
\[
 Q_0S(1\otimes e_0)=0,\qquad
 SQ_0(1\otimes e_0)=e^{-c}(1\otimes e_1),
\]
so $Q_0\notin\{S\}'$.

The same projection commutes with every scalar multiplier, but is not
itself a scalar multiplier. If $Q_0=M_g$, its action on $1\otimes e_0$
would give $g=1$ almost everywhere, whereas its action on
$1\otimes e_1$ is zero. The abelian self-adjoint algebra
$\{M_g:g\in L^\infty(X)\}$ is therefore not maximal abelian in
$\B(H_0)$. Arveson's maximal-abelian inclusion theorem
\cite[Theorem~3.3]{Arveson1967} does not apply to this subalgebra of the
transitive commutant.

\subsection{Atomic finite families and a diffuse joint algebra}
We finish with a spectral property of the constructed weights. This
subsection is independent of the transitivity proof.

Let $|S|=(S^*S)^{1/2}$. Identify $H_0$ with
$L^2(X\times\Z,\mu\times\#)$, where $\#$ is counting measure,
and let $Ue_j=e_{j+1}$ on every fiber. Then
\[
 S=U|S|,\qquad (|S|f)(x,j)=\beta_j(x)f(x,j).
\]
All weights are positive, so $|S|$ is injective with dense range and
$U$ is the unitary polar part of $S$. Define
\[
 A_n=U^{-n}|S|U^n\quad(n\in\Z);
 \qquad (A_nf)(x,j)=\beta_{j+n}(x)f(x,j).
\]
For a family of operators, $W^*$ denotes the smallest strongly closed
unital self-adjoint algebra containing that family. A nonzero projection
in an abelian von Neumann algebra is minimal if it contains no smaller
nonzero projection in that algebra. Such an algebra is \emph{atomic}
when its identity is the strong sum of its minimal projections, and
\emph{diffuse} when it has no minimal projection. We will show that
$W^*(A_n:n\in F)$ is atomic for every finite $F\subset\Z$, whereas
$W^*(A_n:n\in\Z)$ is diffuse.

For any finite $F\subset\Z$ and fixed $j$, the tuple
$(\beta_{j+n}(x))_{n\in F}$ depends on finitely many residues $T_m(x)$.
Because the residues are nested, it depends on one residue $T_M(x)$
and takes finitely many values. Allowing the countably many values of
$j$ gives a countable range for the joint tuple on $X\times\Z$.
Its nonempty level sets have positive measure: on every fixed coordinate
slice they are finite unions of residue classes. These level sets give
the minimal projections of the generated abelian algebra, and their
sum is the identity. Thus $W^*(A_n:n\in F)$ is atomic. In particular,
$|S|$ has a pure-point spectral resolution.

For the full family, the two equal central weights determine the
coordinate slice. They are the only weights equal to $e^{-c}$, since
$a_k^\pm\ge c(k+1)>c$ for every $k\ge1$. At $(x,j)$ the tuple indices
$n$ for which $\beta_{j+n}(x)=e^{-c}$ are therefore exactly
$-j-1$ and $-j$. For every $r\in\Z$, the projection onto the slice
$j=r$, denoted by $P_r$, belongs to $W^*(A_n:n\in\Z)$: it is
\[
 \mathbf1_{\{e^{-c}\}}(A_{-r-1})\,
 \mathbf1_{\{e^{-c}\}}(A_{-r}).
\]
Next fix $m\ge1$ and choose $k_m=\ell_{m-1}\in I_m$. On the slice
$j=r$, the multiplier $A_{k_m-r}$ supplies $\beta_{k_m}(x)$.
For $0\le t<N_m$, put
\[
 \lambda_{m,t}=\exp[-c(k_m+1)-J_m-Dt].
\]
Equation~\eqref{eq:a-plus} shows that the multiplication projection
of the residue rectangle $\{T_m=t\}\times\{r\}$ is
\[
 P_r\,\mathbf1_{\{\lambda_{m,t}\}}(A_{k_m-r}).
\]
Thus every residue-class projection on every coordinate slice belongs
to the generated algebra. Such rectangles generate the product
sigma-algebra. Conversely all $A_n$ are multiplication operators, so,
modulo completion and null sets,
\[
 W^*(A_n:n\in\Z)
   =L^\infty(X\times\Z,\mu\times\#)
\]
in its multiplication representation. Haar measure on $\Z_2$ has no
atoms, and consequently this algebra is diffuse.

The pure-point assertion does not assert discreteness of the spectrum
or compactness of $S$. Nor is this full multiplier algebra the scalar
base algebra used in the transitivity proof. In fact $|S|$ does not
commute with $S$, because the adjacent weights $\beta_0=e^{-c}$ and
$\beta_1\le e^{-2c}$ differ. The polar-conjugate algebra is therefore
not contained in $\{S\}'$.
